\section{ANTECEDENTES}

Los sistemas de votación electrónica han evolucionado desde modelos centralizados, que dependían de una única autoridad, hacia arquitecturas basadas en cadenas de bloques (blockchain). Esta descentralización dificulta que se alteren los resultados de forma retroactiva. Sin embargo, como demostraron investigaciones recientes, aún existen desafíos importantes para garantizar la privacidad del votante y proteger el sistema contra la futura amenaza de las computadoras cuánticas. (Poudel et al., 2024)

En el ámbito de la privacidad y transparencia, el proyecto ``Ordinos'' sentó las bases para que un sistema oculte los resultados parciales durante la elección y permita a cada votante verificar que su voto fue contado correctamente (verificabilidad de extremo a extremo). No obstante, su seguridad matemática es vulnerable a ataques cuánticos. Para solucionar esto, se ha propuesto adaptar estos sistemas utilizando nuevos modelos matemáticos resistentes, como la criptografía basada en retículos. (Wabartha, 2025) Asimismo, para equilibrar la privacidad del usuario con la posibilidad de auditar los votos, las pruebas matemáticas de conocimiento cero se han posicionado como la herramienta más adecuada. (Finogina y Herranz, 2023)

Otras investigaciones han explorado diferentes caminos para asegurar el voto post-cuántico, como el uso de redes que mezclan los votos de forma anónima o esquemas de cifrado verificable. (Farzaliyev, Willemson y Kaasik, 2021; Silde, 2021; Roux-Langlois, 2021) A nivel global, el Instituto Nacional de Estándares y Tecnología (NIST) ha publicado oficialmente los estándares de seguridad post-cuántica. (NIST, 2024) Diferentes estudios de rendimiento han confirmado que estos nuevos esquemas son eficientes y pueden ser ejecutados sin problemas en computadoras y dispositivos comunes. (Lopez, Cadena y Rahman, 2025; Chaurasia, Sharma y Sharma, 2025)

La integración de esta nueva criptografía con pruebas de conocimiento cero y tecnología blockchain es considerada la vía más prometedora para la votación electrónica. (Ezeh et al., 2025; Yu, 2025) Sin embargo, las propuestas actuales son en su mayoría teóricas y carecen de pruebas de rendimiento sobre equipos reales, una brecha que la presente investigación busca abordar.

En el contexto institucional,  la Universidad Mayor de San Andrés UMSA cuenta con investigaciones vinculadas a esta temática. Trabajos pioneros han abordado la encriptación híbrida con distribución cuántica de claves. (Balboa Santander, 2006) De igual forma, se han realizado análisis comparativos entre algoritmos cuánticos y clásicos. (Nina Chura, 2014) La presente propuesta se inserta en esta linea de  investigación académica, con el objetivo de avanzar hacia una solución integral y verificable que proteja los procesos electorales contra las amenazas del futuro.

